\SourceLeaf{084}
subiiciemus, ubi solae aequationes secundae classis respiciuntur. Ita alternis vicibus, modo priorem classem modo posteriorem respicientes, compensationes continuo decrescentes obtinebimus, et si distributio scite adornata fuerat, post paucas iterationes ad numeros stabiles perveniemus. Si plures quam duae classes factae sunt, res simili modo se habebit: classes singulae deinceps in computum venient, post ultimam iterum prima et sic porro. Sed sufficiat hoc loco, hunc modum addigitavisse, cuius efficacia multum utique a scita applicatione pendebit.

\Article{21}
Restat, ut suppleamus demonstrationem lemmatis in art. 8 suppositi, ubi tamen perspicuitatis caussa alias denotationes huic negotio magis adaptatas adhibebimus.

Sint itaque \(x^0,x',x'',x'''\) etc. indeterminatae, supponamusque, ex aequationibus
\[
\begin{aligned}
n^{00}x^0+n^{01}x'+n^{02}x''+n^{03}x'''+\text{etc.}&=X^0\\
n^{10}x^0+n^{11}x'+n^{12}x''+n^{13}x'''+\text{etc.}&=X'\\
n^{20}x^0+n^{21}x'+n^{22}x''+n^{23}x'''+\text{etc.}&=X''\\
n^{30}x^0+n^{31}x'+n^{32}x''+n^{33}x'''+\text{etc.}&=X'''\\
\text{etc.}&
\end{aligned}
\]
sequi per eliminationem has
\[
\begin{aligned}
N^{00}X^0+N^{01}X'+N^{02}X''+N^{03}X'''+\text{etc.}&=x^0\\
N^{10}X^0+N^{11}X'+N^{12}X''+N^{13}X'''+\text{etc.}&=x'\\
N^{20}X^0+N^{21}X'+N^{22}X''+N^{23}X'''+\text{etc.}&=x''\\
N^{30}X^0+N^{31}X'+N^{32}X''+N^{33}X'''+\text{etc.}&=x'''\\
\text{etc.}&
\end{aligned}
\]

Substitutis itaque in aequatione prima et secunda secundi systematis valoribus quantitatum \(X,X',X'',X'''\) etc. e primo systemate, obtinemus
\[
\begin{aligned}
x^0={}&N^{00}(n^{00}x^0+n^{01}x'+n^{02}x''+n^{03}x'''+\text{etc.})\\
&+N^{01}(n^{10}x^0+n^{11}x'+n^{12}x''+n^{13}x'''+\text{etc.})\\
&+N^{02}(n^{20}x^0+n^{21}x'+n^{22}x''+n^{23}x'''+\text{etc.})\\
&+N^{03}(n^{30}x^0+n^{31}x'+n^{32}x''+n^{33}x'''+\text{etc.})\quad\text{etc.}
\end{aligned}
\]
\SourceLeaf{085}
nec non
\[
\begin{aligned}
x'={}&N^{10}(n^{00}x^0+n^{01}x'+n^{02}x''+n^{03}x'''+\text{etc.})\\
&+N^{11}(n^{10}x^0+n^{11}x'+n^{12}x''+n^{13}x'''+\text{etc.})\\
&+N^{12}(n^{20}x^0+n^{21}x'+n^{22}x''+n^{23}x'''+\text{etc.})\\
&+N^{13}(n^{30}x^0+n^{31}x'+n^{32}x''+n^{33}x'''+\text{etc.})\quad\text{etc.}
\end{aligned}
\]

Quum utraque aequatio manifesto esse debeat aequatio identica, tum in priore tum in posteriore pro \(x^0,x',x'',x'''\) etc. valores quoslibet determinatos substituere licebit. Substituamus in priore
\[
x^0=N^{10},\quad x'=N^{11},\quad x''=N^{12},\quad x'''=N^{13}\quad\text{etc.}
\]
in posteriore vero
\[
x^0=N^{00},\quad x'=N^{01},\quad x''=N^{02},\quad x'''=N^{03}\quad\text{etc.}
\]
His ita factis subtractio producit
\[
\begin{aligned}
N^{10}-N^{01}={}&(N^{00}N^{11}-N^{10}N^{01})(n^{01}-n^{10})\\
&+(N^{00}N^{12}-N^{10}N^{02})(n^{02}-n^{20})\\
&+(N^{00}N^{13}-N^{10}N^{03})(n^{03}-n^{30})\\
&+\text{etc.}\\
&+(N^{01}N^{12}-N^{11}N^{02})(n^{12}-n^{21})\\
&+(N^{01}N^{13}-N^{11}N^{03})(n^{13}-n^{31})\\
&+\text{etc.}\\
&+(N^{02}N^{13}-N^{12}N^{03})(n^{23}-n^{32})\\
&+\text{etc. etc.}
\end{aligned}
\]
quae aequatio ita quoque exhiberi potest
\[
N^{10}-N^{01}=\Sigma(N^{0\alpha}N^{1\beta}-N^{1\alpha}N^{0\beta})(n^{\alpha\beta}-n^{\beta\alpha})
\]
denotantibus \(\alpha\beta\) omnes combinationes indicum inaequalium.

Hinc colligitur, si fuerit
\[
\begin{aligned}
&n^{01}=n^{10},\quad n^{02}=n^{20},\quad n^{03}=n^{30}\\
&n^{12}=n^{21},\quad n^{13}=n^{31},\quad n^{23}=n^{32},\quad\text{etc.}
\end{aligned}
\]
sive generaliter \(n^{\alpha\beta}=n^{\beta\alpha}\), fore etiam
\[
N^{10}=N^{01}
\]
\SourceLeaf{086}
Et quum ordo indeterminatarum in aequationibus propositis sit arbitrarius, manifesto in illa suppositione erit generaliter
\[
N^{\alpha\beta}=N^{\beta\alpha}
\]

\Article{22}
Quum methodus in hac commentatione exposita applicationem imprimis frequentem et commodam inveniat in calculis ad geodaesiam sublimiorem pertinentibus, lectoribus gratam fore speramus illustrationem praeceptorum per nonnulla exempla hinc desumta.

Aequationes conditionales inter angulos systematis triangulorum e triplici potissimum fonte sunt petendae.

I. Aggregatum angulorum horizontalium, qui circa eundem verticem gyrum integrum horizontis complent, aequare debet quatuor rectos.

II. Summa trium angulorum in quovis triangulo quantitati datae aequalis est, quum, quoties triangulum est in superficie curva, excessum illius summae supra duos rectos tam accurate computare liceat, ut pro absolute exacto haberi possit.

III. Fons tertius est ratio laterum in triangulis catenam clausam formantibus. Scilicet si series triangulorum ita nexa est, ut secundum triangulum habeat latus unum \(a\) commune cum triangulo primo, aliud \(b\) cum tertio; perinde quartum triangulum cum tertio habeat latus commune \(c\), cum quinto latus commune \(d\), et sic porro usque ad ultimum triangulum, cui cum praecedenti latus commune sit \(k\), et cum triangulo primo rursus latus \(l\), valores quotientium \(\frac al,\allowbreak\frac ba,\allowbreak\frac cb,\allowbreak\frac dc\ .\ .\ .\ .\ \frac lk\), innotescent resp. e binis angulis triangulorum successivorum, lateribus communibus oppositis, per methodos notas, unde quum productum illarum fractionum fieri debeat \(=1\), prodibit aequatio conditionalis inter sinus illorum angulorum, (parte tertia excessus sphaerici vel sphaeroidici, si triangula sunt in superficie curva, resp. diminutorum).

Ceterum in systematibus triangulorum complicatioribus saepissime accidit, ut aequationes conditionales tum secundi tum tertii generis plures se offerant, quam retinere fas est, quoniam pars earum in reliquis iam contenta est. Contra rarior erit casus, ubi aequationibus conditionalibus secundi generis adiungere oportet aequationes similes ad figuras plurium laterum spectantes, puta tunc tantum; ubi
